\section*{Hoofdstuk \ref{ch:b3:colloids} --- Grensvlakken, oppervlakteactieve stoffen en colloïden}

\begin{solution}{exo:b3:colloids:1}
$r = \qty{0.50}{\micro m}$: $\Delta p = 2 \times 0.0720/\num{5.0e-7} = \qty{2.9e5}{Pa}$; binnenin $1.0 + 2.9 = \qty{3.9}{bar}$.
\end{solution}

\begin{solution}{exo:b3:colloids:2}
$\ln(p/p^*) = 2 \times 0.0720 \times \num{1.807e-5}/(\num{1.0e-8} \times 8.314 \times 298.15) = 0.105$; $p/p^* = 1.11$.
\end{solution}

\begin{solution}{exo:b3:colloids:3}
$\cos\theta = (40 - 20)/72 = 0.278$, $\theta = \ang{74}$: onder $90^\circ$; de vloeistof \omterm{def:b3:colloids:wetting}{bevochtigt} de
vaste stof gedeeltelijk.
\end{solution}

\begin{solution}{exo:b3:colloids:4}
NaCl: $I = \qty{10}{mol.m^{-3}}$, $\kappa^{-1} = 0.96\sqrt{100/10} = \qty{3.0}{nm}$. \ce{MgSO4}: $I = \qty{40}{mol.m^{-3}}$, \qty{1.5}{nm}.
\end{solution}

\begin{solution}{exo:b3:colloids:5}
$\Gamma = \num{12.6e-3}/(2 \times 8.314 \times 298.15) = \qty{2.54e-6}{mol.m^{-2}}$; oppervlakte $1/(\Gamma N_A) =
\qty{0.65}{nm^2}$ per molecuul.
\end{solution}

\begin{solution}{exo:b3:colloids:6}
Onder de CMC is de helling $-12.9$ \unit{mN/m} per eenheid van $\log c$ (van $-5.0$ tot $-4.0$); het
plateau ligt bij 33.4. De rechte bereikt het bij $\log c = -4.0 + (39.1 - 33.4)/12.9 = -3.56$: CMC
$\approx \qty{2.8e-4}{mol/L}$.
\end{solution}

\begin{solution}{exo:b3:colloids:7}
SDS: $P = 0.350/(0.60 \times 1.67) = 0.35$, op de grens tussen bollen en korte cilinders:
bolvormige \omterm{def:b3:colloids:micelle}{micellen} in de buurt van de CMC. Lipide: $P = 0.700/(0.65 \times 1.67) = 0.64$:
bilagen, dus vesikels.
\end{solution}

\begin{solution}{exo:b3:colloids:8}
\ce{CaCl2}: $60/64 = \qty{0.94}{mM}$ \ce{Ca^{2+}}; \ce{AlCl3}: $60/729 = \qty{0.082}{mM}$ \ce{Al^{3+}}. \ce{Na3PO4}
coaguleert via zijn \ce{Na+} (bij ongeveer \qty{20}{mM} zout); het fosfaat,
een co-ion, kan zelfs adsorberen en de negatieve lading verhogen.
\end{solution}

\begin{solution}{exo:b3:colloids:9}
Hydrofiel deel $6 \times 44.0 + 17.0 = \qty{281.0}{g/mol}$ van \qty{450.0}{g/mol} (atoommassa's van het boek):
HLB $= 20 \times 281.0/450.0 = 12.5$, een emulgator voor olie-in-water.
\end{solution}

\begin{solution}{exo:b3:colloids:10}
De olie van de kleine druppels is beter oplosbaar in de continue fase (met
de factor $\exp(2\gamma V_{\mathrm m}/rRT)$, waarvan de exponent tien keer groter is voor \qty{50}{nm} dan voor \qty{500}{nm}): ze
diffundeert naar de grote druppels, die groeien terwijl de kleine
verdwijnen. Deze \omterm{def:b3:colloids:ripening}{ostwaldrijping} maakt een \omterm{def:b3:colloids:colloid}{emulsie} grover, zelfs als er
nooit twee druppels versmelten.
\end{solution}

\begin{solution}{exo:b3:colloids:11}
De meniscus is een bolkap met straal $r/\cos\theta$: de vloeistof net eronder staat onder een
druk die $2\gamma\cos\theta/r$ lager is, en de kolom stijgt tot $\rho gh$ dit compenseert. Voor water is
$h = 2 \times 0.072/(997 \times 9.81 \times \num{1.0e-4}) = \qty{0.15}{m}$.
\end{solution}

\begin{solution}{exo:b3:colloids:12}
Bij \qty{100}{mM} is de \omterm{def:b3:electrode-kinetics:double-layer}{debyelengte} kleiner dan \qty{1}{nm} en sterft de
afstoting uit vóór de
aantrekking: er blijft geen barrière over (voorbij de kromme van \qty{50}{mM}). Een vrij polymeer
werkt niet aan het oppervlak; een laag geënt polymeer stoot bij contact
af, wat het zout ook is, door de entropie die verloren gaat wanneer de
ketens overlappen.
\end{solution}

\begin{solution}{pb:b3:colloids:1}
\textbf{1.} Substituties in het kristalrooster (\ce{Al^{3+}} voor \ce{Si^{4+}}, \ce{Mg^{2+}} voor \ce{Al^{3+}}) laten
een permanente negatieve lading achter, gecompenseerd door uitwisselbare
kationen.
\textbf{2.} $I = \frac12(1.0 + 1.0) = \qty{1.0}{mM}$.
\textbf{3.} $\kappa^{-1} = \qty{9.6}{nm}$.
\textbf{4.} $a\kappa = 10$: de dubbellaag is dun vergeleken met het deeltje,
de voorwaarde voor de benadering van Derjaguin.
\textbf{5.} $v = 2 \times 1650 \times 9.81 \times (\num{1.0e-7})^2/(9 \times \num{0.89e-3}) = \qty{4.0e-8}{m/s}$, ongeveer \qty{3.5}{mm} per dag.
\textbf{6.} Hun diffuse lagen stoten elkaar af voordat de
vanderwaalsaantrekking de overhand neemt.
\textbf{7.} $V = 2\pi\varepsilon_r\varepsilon_0a\psi^2\ln(1 + \eu^{-\kappa h}) - Aa/12h$.
\textbf{8.} Ongeveer $39\,kT$.
\textbf{9.} Bij een fractie van de orde $\eu^{-39} \approx 10^{-17}$ van de botsingen: praktisch nooit.
\textbf{10.} Ze verdwijnt: elke botsing leidt tot contact.
\textbf{11.} Een kortere \omterm{def:b3:electrode-kinetics:double-layer}{debyelengte} beperkt de afstoting tot kleine
spleten, waar de
aantrekking, die groeit als $1/h$, overheerst.
\textbf{12.} De CCC varieert als $z^{-6}$ van de lading van het tegenion.
\textbf{13.} $50/64 = \qty{0.78}{mM}$.
\textbf{14.} $50/729 = \qty{0.069}{mM}$.
\textbf{15.} De kationen zijn de tegenionen van de negatieve deeltjes: ze
hopen zich op in de dubbellaag en schermen haar af.
\textbf{16.} Sulfaat is een co-ion en wordt door de deeltjes afgestoten.
\textbf{17.} \qty{0.069}{mol} \ce{Al^{3+}} per kubieke meter.
\textbf{18.} \qty{0.034}{mol} $\ce{Al2(SO4)3}$.
\textbf{19.} $0.0343 \times 342.3 = \qty{11.7}{g}$.
\textbf{20.} \qty{11.7}{kg} per uur.
\textbf{21.} $\ce{Al^3+ + 3 HCO3- -> Al(OH)3 + 3 CO2}$.
\textbf{22.} $3 \times 0.069 = \qty{0.21}{mM}$ van de \qty{1.0}{mM}: een vijfde. Het resterende
waterstofcarbonaat en het opgeloste \ce{CO2} bufferen het water: de pH daalt
maar weinig.
\textbf{23.} De sterk geladen aluminiumdeeltjes adsorberen en keren de
lading van de deeltjes om, die elkaar dan opnieuw afstoten.
\textbf{24.} \textbf{Ongeveer \qty{12}{g} aluminiumsulfaat per kubieke meter} (\qty{11.7}{g}), vóór de extra
hoeveelheid die in de praktijk voor de hydrolyse nodig is.
\end{solution}
